\documentclass[12pt]{article}

\usepackage{sbc-template}
\usepackage{graphicx,url}
\usepackage[utf8]{inputenc}
\usepackage[brazil]{babel}
% ATENÇÃO: o template original da SBC carrega \usepackage[latin1]{inputenc}
% logo após a linha acima. Isso gera erro de compilação ("Option clash for
% package inputenc") porque o pacote é carregado duas vezes com opções
% diferentes. Como este arquivo é salvo em UTF-8, mantenha SOMENTE a linha
% [utf8]{inputenc} e NÃO adicione a linha [latin1]{inputenc}.

\sloppy

\title{Previsão de Produtividade Municipal da Soja no Rio Grande do Sul com Aprendizado de Máquina e Variáveis Agrometeorológicas}

\author{Isabelle Franco, Gustavo Rodrigues, Pedro Henrique, Lorenzo Tadeo, Ivan Carlos A. de Oliveira\inst{1}}

\address{Faculdade de Computação e Informática (FCI) -- Universidade Presbiteriana
  Mackenzie\\
  São Paulo -- SP -- Brasil
  \email{\{isabelle.franco,gustavo.rodrigues,pedro.henrique,lorenzo.tadeo\}@mackenzista.com.br}
}

\begin{document}

\maketitle

\begin{center}
\footnotesize
Isabelle Franco (RA 10425395) -- Gustavo Rodrigues (RA 10403091) -- Pedro Henrique (RA 10388298) -- Lorenzo Tadeo (RA 10420067)\\
Orientador: Prof.\ Dr.\ Ivan Carlos Alcântara de Oliveira
\end{center}

\begin{abstract}
  Rio Grande do Sul shows substantial year-to-year variability in soybean
  yield, driven by climatic conditions and by fertilization practices that
  vary across municipalities. This project proposes a machine learning
  approach to anticipate municipal soybean yield in the state using public
  data. The dataset combines historical yield, area and production series
  from IBGE/SIDRA Table 1612 (RS municipalities, 2000--2023); the Oceanic
  Niño Index (ONI) from NOAA, as a large-scale climate signal linked to
  regional rainfall patterns; and observed INMET station data, used to
  validate the meteorological variables. As a distinguishing feature relative
  to approaches that rely solely on climate and yield history, the project
  incorporates attributes related to common phosphate fertilization in RS --
  MAP, Simple Superphosphate, Triple Superphosphate and DAP -- based on
  agronomic recommendations from Embrapa and the RS/SC Soil Fertility
  Committee, aiming to capture the effect of phosphorus, sulfur and calcium
  availability on yield. In N1, the work focuses on problem definition,
  ethical discussion, dataset documentation, exploratory analysis and data
  preparation in Python, along with the methodological design. In later
  stages, Linear Regression, Random Forest and Gradient Boosting will be
  compared against historical baselines under temporal validation, with SHAP
  applied to the best model to identify the climatic and fertilization
  factors most associated with yield variation.
\end{abstract}

\begin{resumo}
  O Rio Grande do Sul apresenta elevada variabilidade interanual na
  produtividade da soja, influenciada por condições climáticas e por
  práticas de adubação que variam entre municípios. Este projeto propõe uma
  solução de Inteligência Artificial para estimar antecipadamente a
  produtividade municipal da soja no estado, a partir de dados públicos. A
  base será estruturada com a série histórica de rendimento, área e produção
  da Tabela 1612 do IBGE/SIDRA (municípios do RS, 2000--2023); o Oceanic
  Niño Index (ONI) da NOAA, como sinal climático de grande escala associado
  ao regime de chuvas na região; e dados observados do INMET, utilizados
  para validação das variáveis meteorológicas. Como diferencial em relação a
  abordagens que consideram apenas clima e histórico produtivo, serão
  incorporados atributos relacionados à adubação fosfatada usual no RS --
  MAP, Superfosfato Simples, Superfosfato Triplo e DAP --, com base em
  recomendações agronômicas da Embrapa e da Comissão de Química e
  Fertilidade do Solo RS/SC, buscando capturar o efeito da disponibilidade
  de fósforo, enxofre e cálcio no solo gaúcho sobre a produtividade. Na N1, o
  trabalho concentra-se na definição do problema, discussão ética,
  documentação do dataset, análise exploratória e preparação dos dados em
  Python, além da especificação metodológica. Nas etapas seguintes, serão
  comparados modelos de Regressão Linear, Random Forest e Gradient Boosting
  com linhas de base históricas, sob validação temporal, e aplicada a
  técnica SHAP ao melhor modelo para identificar os fatores climáticos e de
  adubação mais associados às oscilações de produtividade.
\end{resumo}

\section{Introdução}

\subsection{Contextualização}

A soja é uma das principais commodities do agronegócio brasileiro, e o Rio
Grande do Sul é um dos estados mais expressivos nessa cultura. A
produtividade varia significativamente entre municípios e safras,
condicionada pela interação entre regime de chuvas, temperatura, ocorrência
de veranicos, fenômenos climáticos de larga escala (El Niño/La Niña) e
práticas de manejo do solo -- em especial a adubação fosfatada, relevante
para os solos ácidos e historicamente pobres em fósforo do RS. A Pesquisa
Agrícola Municipal (PAM) do IBGE, disponibilizada pela Tabela 1612 do SIDRA,
oferece uma série histórica municipal de rendimento, área e produção desde
1974, o que permite trabalhar com uma amostra ampla (muitos municípios
$\times$ muitas safras) mantendo uma única granularidade geográfica -- o
município -- dentro de um único estado.

\subsection{Justificativa}

Modelos de Machine Learning como Random Forest e Gradient Boosting são
adequados a este problema porque capturam relações não lineares entre
variáveis climáticas, de adubação e produtividade, superando as limitações
de projeções estatísticas lineares tradicionais. Complementar o histórico
produtivo com o ONI (sinal climático disponível antes e durante a safra) e
com atributos de adubação fosfatada permite testar se informações
conhecidas \emph{antes} do fim da safra carregam sinal preditivo útil -- e
não apenas reproduzir, a posteriori, o que já é sabido pela própria série de
rendimento.

\subsection{Motivação}

Estimativas antecipadas e municipais de produtividade podem apoiar o
planejamento de cooperativas, produtores e órgãos públicos, sinalizando com
maior granularidade geográfica municípios sob risco de rendimento abaixo da
referência histórica. O projeto não pretende substituir avaliações
agronômicas locais nem prever o resultado de uma propriedade específica -- o
objetivo é um indicador regional, transparente e reprodutível, construído
inteiramente a partir de dados públicos.

\subsection{Objetivo}

Desenvolver e avaliar uma abordagem de aprendizado de máquina supervisionado
capaz de prever a produtividade da soja (kg/ha) por \textbf{município do Rio
Grande do Sul e por safra}, a partir de variáveis climáticas, do sinal ENOS
(ONI) e de atributos de adubação fosfatada, comparando o desempenho dos
modelos contra linhas de base históricas simples.

\subsection{Contribuições esperadas}

\begin{itemize}
  \item Consolidar uma base analítica integrada em nível
    município--safra, cobrindo os municípios do RS e as safras disponíveis
    na Tabela 1612 do IBGE/SIDRA (2000--2023);
  \item Construir atributos climáticos agregados por safra a partir do ONI e
    validados com dados observados do INMET;
  \item Incorporar atributos de adubação fosfatada (MAP, SS, ST, DAP) com
    base em recomendações agronômicas por tipo de solo;
  \item Treinar e comparar Regressão Linear (baseline), Random Forest e um
    modelo de Gradient Boosting contra linhas de base históricas (média
    histórica do município e rendimento da safra anterior), sob validação
    temporal;
  \item Aplicar SHAP ao melhor modelo para identificar os fatores climáticos
    e de adubação mais influentes sobre a produtividade.
\end{itemize}

\section{Fundamentação Teórica (Resumida)}

\subsection{Conceitos fundamentais}

O problema é formulado como \textbf{regressão supervisionada}, em que a
variável-alvo é o rendimento médio municipal da soja (kg/ha). Serão
comparados três modelos: \textbf{Regressão Linear}, usada como baseline
interpretável; \textbf{Random Forest}, um ensemble de árvores de decisão por
\emph{bagging}, adequado para capturar interações não lineares entre
variáveis climáticas e de solo; e \textbf{Gradient Boosting}, que constrói
árvores sequencialmente para corrigir os erros residuais do conjunto
anterior, geralmente com melhor desempenho em dados tabulares. Para
interpretar as previsões do melhor modelo será usado \textbf{SHAP} (SHapley
Additive exPlanations), técnica de Inteligência Artificial Explicável (XAI)
que atribui a cada variável uma contribuição marginal para a previsão
individual, permitindo identificar quais fatores climáticos e de adubação
mais pesam nas estimativas.

\subsection{Trabalhos relacionados}

Diversos estudos aplicam aprendizado de máquina à previsão de produtividade
agrícola. Jeong et al. (2016) mostraram que Random Forest supera modelos de
regressão múltipla na previsão de produtividade de grãos em múltiplas
regiões do mundo. Van Klompenburg, Kassahun e Catal (2020), em revisão
sistemática, identificam Random Forest e redes neurais como os algoritmos
mais utilizados para previsão de safra, destacando a importância de
variáveis climáticas e a necessidade de validação temporal para evitar
vazamento de dados. Especificamente para a soja no Sul do Brasil, Pilau et
al. (2022) e Nóia Júnior et al. (2020) documentam efeito negativo
consistente da fase fria do ENOS (La Niña) sobre a produtividade da soja no
Rio Grande do Sul, o que justifica a inclusão do ONI como atributo. Von Bloh
et al. (2023) comparam modelos de machine learning para previsão de
produtividade de soja no Brasil e reforçam a vantagem de modelos baseados em
árvores frente a baselines lineares quando há relações não lineares entre
clima e rendimento. Nenhum desses trabalhos, no entanto, incorpora
explicitamente variáveis de adubação fosfatada municipal -- o que constitui
a principal diferença da proposta deste projeto.

\section{Metodologia e Resultados Esperados}

\subsection{Abordagem}

A pesquisa é aplicada, quantitativa e preditiva, com unidade de análise
município--safra, restrita ao Rio Grande do Sul. Os dados serão divididos
\textbf{temporalmente} -- treino com safras mais antigas, teste com as
safras mais recentes --, nunca por divisão aleatória, para evitar
vazamento de informação futura. Os hiperparâmetros dos modelos serão
ajustados apenas sobre os dados de treino, com validação temporal interna
(ex.: \emph{walk-forward}).

Os três modelos (Regressão Linear, Random Forest, Gradient Boosting) serão
comparados entre si e contra duas linhas de base simples: a média histórica
expandida do município (calculada apenas com safras anteriores à safra
prevista) e o rendimento da safra imediatamente anterior. As métricas de
avaliação serão MAE, RMSE e R\textsuperscript{2}, reportadas de forma
agregada e, quando a amostra permitir, também por município.

\subsubsection{Justificativa: regressão em painel com validação temporal, e
  não série temporal clássica}

Optou-se por tratar o problema como uma \textbf{regressão em painel}
(municípios empilhados, com o rendimento da safra anterior e a média
histórica expandida como atributos) sob \textbf{validação temporal}, em vez
de um modelo de série temporal clássico por município (ARIMA, suavização
exponencial, LSTM univariado). A razão é o volume de dados disponível por
unidade: a série de cada município cobre cerca de 24 safras (2000--2023),
abaixo do mínimo prático ($\sim$40--50 pontos) para ajustar um ARIMA ou uma
rede recorrente de forma confiável; no RS, isso é agravado pela criação de
novos municípios por desmembramento ao longo do período, que interrompe a
continuidade de algumas séries.

Duas alternativas de série temporal ``de verdade'' foram consideradas e
descartadas para esta etapa:

\begin{enumerate}
  \item \textbf{Série temporal agregada por estado} (produtividade média do
    RS por safra, $\sim$50 pontos desde 1974, viável para ARIMA) --
    descartada por eliminar a granularidade municipal, contrariando
    diretamente a exigência do parecer de fixar uma granularidade única e
    informativa.
  \item \textbf{Modelo de série temporal global} (ex.: DeepAR, N-BEATS,
    Temporal Fusion Transformer), treinado simultaneamente sobre todas as
    séries municipais -- reduz a exigência de pontos por série, mas
    tipicamente ainda pede centenas de séries com dezenas de pontos cada,
    além de maior complexidade de implementação e ajuste; descartada por
    desproporção entre esforço e prazo do N1/N2 (três meses).
\end{enumerate}

A abordagem de painel preserva a estrutura temporal essencial -- o treino
usa apenas safras passadas para prever safras futuras (split temporal, nunca
aleatório) -- sem exigir o volume de dados que os métodos de série temporal
genuínos demandariam. Essa escolha poderá ser revisitada como extensão,
caso o tempo do projeto permita testar um modelo de série temporal agregado
por estado como comparação adicional.

\subsection{Dataset: origem, conteúdo e preparação}

\begin{table}[ht]
\centering
\caption{Fontes de dados utilizadas no projeto}
\label{tab:dataset}
\begin{tabular}{|p{2.3cm}|p{4.7cm}|p{4.2cm}|}
\hline
\textbf{Fonte} & \textbf{Conteúdo} & \textbf{Uso no projeto} \\ \hline
IBGE/SIDRA -- Tabela 1612 & Rendimento (kg/ha), área plantada/colhida e
  produção de soja, por município do RS, 2000--2023 & Variável-alvo e
  histórico produtivo \\ \hline
NOAA/CPC -- ONI & Índice trimestral do ENOS (El Niño/La Niña) & Sinal
  climático de grande escala, disponível antes e durante a safra \\ \hline
INMET & Precipitação e temperatura observadas em estações & Validação da
  qualidade/cobertura dos dados climáticos usados \\ \hline
Embrapa / CQFS-RS/SC & Recomendações de adubação fosfatada (MAP,
  Superfosfato Simples, Superfosfato Triplo, DAP) por classe de solo &
  Atributos de disponibilidade de P, S e Ca no solo \\ \hline
\end{tabular}
\end{table}

A base final terá como chaves código IBGE do município e safra. Os dados
brutos do SIDRA serão obtidos via API do IBGE (consultando primeiro os
metadados da Tabela 1612 para confirmar os códigos exatos de variável e
classificação antes da extração) e reorganizados para o formato
município $\times$ ano; o ONI será agregado por janela de safra
(semeadura--colheita); e os dados do INMET, obtidos via BDMEP, usarão a
estação mais próxima de cada município como referência de validação
cruzada com o sinal climático agregado. A análise exploratória cobrirá:
cobertura temporal e espacial (municípios com séries incompletas), valores
ausentes, correlação entre rendimento e as demais variáveis, e a relação
entre fase do ENOS e rendimento médio -- reproduzindo, com a base do RS, o
tipo de verificação já documentado na literatura (Pilau et al., 2022).

\begin{quote}
\textbf{Nota sobre o índice climático:} a NOAA passou a adotar o Índice ONI
Relativo (RONI) como referência oficial para monitoramento e previsão do
ENOS, em substituição ao ONI tradicional. O projeto usará o ONI por ser o
índice mais consolidado na literatura de produtividade agrícola no RS
(Pilau et al., 2022; Nóia Júnior et al., 2020), mas essa escolha será
justificada explicitamente no artigo, com o RONI como alternativa a avaliar
caso o tempo do projeto permita.
\end{quote}

\begin{quote}
\textbf{Limitação a declarar explicitamente:} os dados de adubação
fosfatada normalmente não existem como série real de consumo por município
e ano, mas como recomendação agronômica por tipo de solo. Isso será tratado
como atributo de referência (dose recomendada por classe de solo/região), e
essa limitação será documentada na seção de aspectos éticos e nas
conclusões -- não será apresentada como o volume efetivamente aplicado por
cada produtor.
\end{quote}

\section{Resultados Parciais}

\subsection{Aspectos éticos do uso da IA e responsabilidade}

O projeto utiliza exclusivamente dados públicos e agregados em nível
municipal, sem dados pessoais ou identificação de indivíduos, o que o
mantém aderente à LGPD e dispensa submissão a Comitê de Ética. O principal
risco ético está na interpretação do indicador: ele não deve ser tomado
como recomendação automática de crédito, seguro ou decisão de plantio, e
decisões com impacto econômico devem continuar envolvendo avaliação
agronômica humana. Municípios com séries curtas ou estações meteorológicas
com falhas podem apresentar desempenho inferior do modelo -- essa
limitação será documentada, assim como a natureza normativa (e não
observacional) dos dados de adubação usados. Código, versões das bases e
parâmetros dos modelos serão versionados no GitHub do projeto para
permitir auditoria e reprodução.

\subsection{Resultados parciais obtidos até o momento}

% TODO (grupo): preencher conforme o que já foi executado -- coleta via API
% do SIDRA, primeira limpeza do ONI, integração inicial, gráficos da análise
% exploratória.

\subsection{Endereço GitHub}

% TODO (grupo): inserir o link do repositório público do grupo, por
% exemplo:
\url{https://github.com/USUARIO/NOME-DO-REPOSITORIO}

\begin{thebibliography}{99}

\bibitem{breiman2001}
BREIMAN, L. Random forests. \emph{Machine Learning}, v. 45, p. 5--32, 2001.

\bibitem{friedman2001}
FRIEDMAN, J. H. Greedy function approximation: a gradient boosting machine.
\emph{The Annals of Statistics}, v. 29, n. 5, p. 1189--1232, 2001.

\bibitem{ibge}
IBGE -- INSTITUTO BRASILEIRO DE GEOGRAFIA E ESTATÍSTICA. Produção Agrícola
Municipal: tabela 1612 -- lavouras temporárias. Rio de Janeiro: IBGE/SIDRA.
Disponível em: \url{https://sidra.ibge.gov.br/tabela/1612}. Acesso em: 24
set. 2026.

\bibitem{inmet}
INMET -- INSTITUTO NACIONAL DE METEOROLOGIA. BDMEP -- Banco de Dados
Meteorológicos. Brasília, DF. Disponível em:
\url{https://bdmep.inmet.gov.br/}. Acesso em: 24 set. 2026.

\bibitem{jeong2016}
JEONG, J. H. et al. Random forests for global and regional crop yield
predictions. \emph{PLOS ONE}, v. 11, n. 6, e0156571, 2016.

\bibitem{noaa}
NOAA -- NATIONAL OCEANIC AND ATMOSPHERIC ADMINISTRATION. Climate Prediction
Center. Oceanic Niño Index (ONI): Cold \& Warm Episodes by Season.
Disponível em:
\url{https://origin.cpc.ncep.noaa.gov/products/analysis_monitoring/ensostuff/ONI_v5.php}.
Acesso em: 24 set. 2026.

\bibitem{noia2020}
NÓIA JÚNIOR, R. S. et al. Effects of the El Niño Southern Oscillation
phenomenon and sowing dates on soybean yield and on the occurrence of
extreme weather events in southern Brazil. \emph{Agricultural and Forest
Meteorology}, v. 290, 108038, 2020.

\bibitem{pilau2022}
PILAU, F. G. et al. Impact of ENSO-related rainfall variability on soybean
yield in the state of Rio Grande do Sul, Brazil. \emph{Agrometeoros}, v. 30,
e027115, 2022.

\bibitem{sbcs2016}
SOCIEDADE BRASILEIRA DE CIÊNCIA DO SOLO -- NÚCLEO REGIONAL SUL / COMISSÃO
DE QUÍMICA E FERTILIDADE DO SOLO -- RS/SC. Manual de Calagem e Adubação
para os Estados do Rio Grande do Sul e de Santa Catarina. Porto Alegre,
2016.

\bibitem{vanklompenburg2020}
VAN KLOMPENBURG, T.; KASSAHUN, A.; CATAL, C. Crop yield prediction using
machine learning: a systematic literature review. \emph{Computers and
Electronics in Agriculture}, v. 177, 105709, 2020.

\bibitem{vonbloh2023}
VON BLOH, M. et al. Machine learning for soybean yield forecasting in
Brazil. \emph{Agricultural and Forest Meteorology}, v. 341, 109670, 2023.

\end{thebibliography}

\end{document}
